\section{The Three Readings Compared}
\label{sec:comparison}

The three readings share their ground. In each, one King calls by grace and
judges those he has called; the Church is a mixed company until he comes in,
and the man without the garment has come in under the call. Each hears the
Epistle's commands as the life the called are to lead, and each ends with the
Communion's wish, \latin{utinam dirigantur viae meae}, which St Augustine, St
Hilary and the Greek \work{Expositiones in Psalmos} printed under St
Athanasius's name all read as a prayer for the grace to do what is commanded.
The witnesses are shared too. St Augustine and St John Chrysostom speak in all
three, Augustine for different loci in each; St Gregory the Great carries the
first and the third and gives the second its principal alternative; and St
Hilary's exposition of the Gradual stands in all three, as part of the first
reading's moral sense, as the second's prayer of the last age, and as the
third's other answer to what the evening sacrifice is.

They differ in the question they put to the texts, and so in which part of
the Gospel and which chants weigh most in each.

\begin{comparisontable}{Question}{The garment}{The call to the highways}{The feast that now is}
The question it answers & What must the guest who has come in be wearing? &
Who was called, by whom, and what became of the call when it was refused? &
Where is the feast now, and how does this assembly take its place at it?\\
Where the Gospel weighs most & Verses 11--14: the king's entry, the garment,
the silence, the binding & Verses 1--10: the two sendings, the refusal, the
city, the highways & Verses 10--11: the hall filled with bad and good, and the
king coming in to see\\
The chants that lead & The Communion's command and wish, with the Gradual's
lifted hands & The Introit's psalm of Israel's history and the Alleluia's
command to the Gentiles & The Gradual's evening prayer, the Offertory's walk in
tribulation, and the Communion\\
What the wedding is & The Church, entered by faith and judged by charity &
The Son's union with a Church of Jews and Gentiles, its glory received in the
resurrection (Hilary) & The present Church at the Lord's Table and the banquet
of the Word\\
The Gradual's lifted hands & Hands lifted in the works of mercy (Hilary),
without anger (Chrysostom) & The prayer of the whole Church still on its road
(Augustine), in the last age (Hilary) & The evening sacrifice of the Passion,
in which the members are figured (Augustine)\\
The strongest difficulty & Whether the garment is gift or work, with Hilary's
baptismal gift beside Gregory's and Augustine's charity & Whether the burned
city is Jerusalem or the persecutors in eternal fire, and Gregory's highways,
which name no people & That the table is not the garment (Augustine), and that
Hilary's evening is not Augustine's\\
Carried by & Gregory, Augustine and Jerome & Chrysostom, Irenaeus and Hilary &
Augustine, Gregory and Chrysostom\\
\end{comparisontable}

No doctrine of one reading is denied by another. The differences inside the
readings are real: whether the garment was given in
baptism or is the charity the baptized may lack, whether the city is Jerusalem
or the persecutors' flesh in eternal fire, whether the highways are the
nations, the teachings of the Gentiles or the failure of worldly undertakings,
and whether the evening sacrifice is the Passion or the works of mercy. But
between the three readings the relation is one of place and time rather than
of contradiction. The second supplies the history in which the other two take
place: the guests examined for the garment and the assembly at the table are
the ones gathered from the highways. The first asks of each of those guests
what the king will look for when he comes in, and the third sets the same
guests at the table now, before he comes.

Each accounts for something the others leave aside. The first alone gives full
weight to the Epistle, reading its four acts as the garment's visible side and
its new man as Christ put on; with St Jerome and St Thomas Aquinas it hears
Christ named in both lessons. The second alone gives the Introit's psalm and the
Alleluia's command their own voice, the psalm of a people told its history and
the command to tell it to the Gentiles, and alone takes up the first half of
the parable, where most of its words are spent. The third alone reads the Mass's
own prayers as prayers of the assembly at the altar: the Gradual that the priest
says again over the altar at a Mass with incense, the Offertory that opens the
offering, the Secret whose words the \latin{Placeat} repeats, and the
Postcommunion's healing.

The Gospel's last sentence shows how they meet. ``For many are called, but few
are chosen.'' In the first reading the saying falls on the guest who came in:
St Gregory's \latin{Quia vocati sumus, novimus; si sumus electi, nescimus}, we
know that we are called, we do not know whether we are chosen. In the second it
falls on the whole history of the call: St Irenaeus's ``as in the former
covenant, `with many of them was He not well pleased;' so also is it the case
here'', and St Hilary's \latin{non \dots\ paucitas in invitatis, sed raritas in
electis}. In the third it falls on the feast that now is, which St Augustine
tells his hearers to receive worthily so that they may attain to the other,
``where no bad men shall come''. One sentence, one judgement, three times: the
guest's, the history's, and the table's.
